\documentclass[conference,compsoc]{IEEEtran}
\IEEEoverridecommandlockouts
\usepackage[utf8]{inputenc}
\usepackage[T1]{fontenc}
\usepackage[nocompress]{cite}
\usepackage{amsmath,amssymb}
\usepackage{graphicx}
\usepackage{booktabs}
\usepackage{tabularx}
\usepackage{array}
\usepackage{url}
\usepackage[hidelinks]{hyperref}
\hyphenation{op-tical net-works semi-conduc-tor}

\begin{document}

\title{Beyond Provenance: Owner-Bound Egress for Data-Exfiltration Defense in Multi-Agent LLM Systems}

\author{\IEEEauthorblockN{Anonymous Author(s)}
\IEEEauthorblockA{Anonymous submission}}

\maketitle

\begin{abstract}
LLM agents increasingly work in teams, handing tasks to each other as messages. A prompt injection relayed by one agent reaches the next in the same form, as a request from a teammate. We call this the hand-off dilemma and study it for data exfiltration. We show that any enforcement rule that decides from an instruction's provenance (its sender, the sender's privilege, or what content the agent has read) receives identical evidence for a legitimate hand-off and an injected one, so it must either block hand-offs or admit the attack. Identity and permission checks do not help, because the leaking agent is allowed to send. We propose owner-bound egress, which judges the effect of an action instead: agents pass records by reference, and an external gateway releases a record only to a recipient bound to its owner. We evaluate seven gateway rules in simulated ten-agent systems with five open-weight models from four families (700 runs, 600 with an attack; 31{,}500 decisions; each rule applied to every run). Privileged agents that had read an injected instruction emailed the attacker's address in 18.5\% of such decisions when told to accept requests from teammates, against 0.2\% when the operator assigned the task directly. All three provenance-based rules completed none of the 700 hand-off tasks. Owner-bound egress blocked all 250 harmful sends and completed 673 of 700 hand-off tasks, the same as with no defence. Its protection has a limit: an agent that copies data inline instead of passing it by reference escapes the binding.
\end{abstract}

\IEEEpeerreviewmaketitle

\section{Introduction}\label{sec:1}
The paradigm of large language model (LLM) utilization has shifted from passive information retrieval to autonomous, agentic execution. Modern LLM agents actively interface with digital environments by parsing communications and source code, invoking external tools, and routing data between third-party services. Consequently, deployment strategies are shifting toward multi-agent systems in which an orchestrating agent decomposes high-level objectives and delegates sub-tasks to subordinate agents using natural language. The standardization of these communication channels and tool integrations is currently driven by emerging protocols like A2A~\cite{a2a2025} and MCP~\cite{mcp2024}. The industry's reliance on these standards is underscored by A2A's designation as a Linux Foundation project, with founding members including AWS, Cisco, Google, Microsoft, Salesforce, SAP, and ServiceNow~\cite{googlecloud2025a2alf}. Nevertheless, as agentic systems acquire deeper integration with sensitive corporate data and external channels, the potential liabilities and cascading consequences of executing compromised instructions increase substantially.

The most direct route to such an action is prompt injection, which OWASP ranks first among risks to LLM applications~\cite{owasp2025llm01}. An instruction planted in a web page, file or message is read by an agent as if it came from the user~\cite{greshake2023}, and in a multi-agent system it can be relayed from one agent to the next~\cite{lee2024promptinfection,triedman2025}. This has already happened in shipped products. In November 2025, a low-privilege user of ServiceNow Now Assist planted text in a ticket that led one agent to recruit a second agent with an administrator's privileges, which copied a ticket the attacker could not read into one they could, with injection protection enabled~\cite{costello2025appomni}. Single-agent incidents follow the same pattern: a public GitHub issue led a coding agent to publish private repository data~\cite{invariant2025github}, and one email led Microsoft 365 Copilot to send internal data to an attacker's server (CVE-2025-32711)~\cite{reddy2025echoleak}. In each case the action that leaked data was one the agent was permitted to take, and this is what makes such attacks hard to stop.

Existing defences struggle with permitted actions because they decide from who is acting or where an instruction came from. Identity defences address who is acting. Their standard form is a reference monitor~\cite{anderson1972}, a single trusted component outside the agents that checks every action, owns keys, scopes and the audit log, and so removes spoofing, scope over-claims and log tampering. It cannot refuse a relayed injection, because the leaking agent is a confused deputy~\cite{hardy1988,triedman2025,ji2026}: it is who it claims to be and holds the permission it uses. Provenance defences go further by tracking the origin of each instruction and restricting agents that have read untrusted input~\cite{debenedetti2025camel,costa2025}, and recent work carries these labels across agent boundaries~\cite{petersgill2026}. In information-flow terms these are integrity labels: they record where an instruction came from, not who may read the data it moves. All of them rely on a premise, that the origin of an instruction separates a legitimate action from an injected one. The premise holds for a single agent serving a trusted user, but delegation removes it.

When agents may instruct each other, a provenance rule receives the same evidence for a legitimate hand-off and for a relayed injection: the same sender, privilege, channel and kind of content (Figure~\ref{fig:1}). It must therefore give both the same verdict, so it either blocks delegation or admits the attack. We call this the hand-off dilemma. The agents themselves show the same limit. Privileged agents acted on an injected instruction in 2 of 1,088 exposures when the operator assigned their task, but in 210 of 1,137 when told to handle requests passed on by teammates. At the same position in the graph, privileged agents given no operator instruction at all acted in 14.2\% of exposures, so what protected the first group was a task the operator had named, and delegation takes that protection away. Since neither the gateway nor the model can separate the two cases by origin, a defence needs a different input.

\begin{figure}[!t]
\centering
\includegraphics[width=\linewidth]{Figure_1.pdf}
\caption{The hand-off dilemma. A real hand-off $E$ and its attack twin $E'$ reach the privileged agent from the same helper, with the same privilege and channel, and both send email to the same recipient $r$, here Eve's registered contact. Only the record set $D$ differs: Eve's own record in $E$, Alice's record in $E'$. A provenance-determined rule does not read $D$, so it gives both the same verdict and either blocks the hand-off or admits the attack.}
\label{fig:1}
\end{figure}

We propose owner-bound egress, a gateway rule that decides on the effect of an action rather than its origin. Agents hold references to records instead of their contents, and the gateway releases a record only to a recipient bound to that record's owner. Because the rule never reads who asked or what the agent has seen, the dilemma does not apply to it. The release check itself is the confidentiality half of classical information-flow control: in the decentralized label model, data carries its owner's policy on who may read it~\cite{myers1997}, and CaMeL attaches the same kind of readers label to values inside a single agent~\cite{debenedetti2025camel}. What delegation changes is where such a check can hold. Inside one agent it is bypassed as soon as a peer's request arrives as a trusted instruction~\cite{petersgill2026}, so we place it outside every agent, at the point where data leaves the system.

We evaluate the rule against six alternatives in simulated ten-agent systems with five open-weight models from four families (700 runs and 31{,}500 decisions), and against scripted insiders that hold valid credentials. Owner-bound egress held all 250 harmful sends, completed as many hand-off tasks as no defence (673 of 700), and failed, as its precondition predicts, only when an agent copied record contents into the message body. The same logs also show why the agents themselves cannot be the line of defence: in every model, whether a privileged agent obeys an injection depends strongly on whether the operator named its task. This paper makes three contributions.

\begin{itemize}
  \item \textbf{The hand-off dilemma.} We define provenance-determined and effect-determined egress rules and prove that every provenance-determined rule gives a legitimate hand-off and its attack twin the same verdict, so its safety and usefulness on twin pairs sum to exactly one (Section~\ref{sec:4}). This bounds a whole class of defences rather than one design, and shows that tracking where instructions came from, however finely, cannot secure delegation. It turns the design question into a test that any defence can be put through, whether its verdict can change when only the data being moved changes, which we apply to prior defences in Table~\ref{tab:2}.
  \item \textbf{A controlled measurement of delegation risk.} Across 31{,}500 logged decisions from five models, privileged agents acted on injections in 0.2\% of exposures when the operator named their task, 14.2\% with no instruction, and 18.5\% when the task was delegated, while relay stayed below one new relayer per relayer (Sections~\ref{sec:7.1} and~\ref{sec:7.4}). This shows that the protection comes from the operator naming the task, that delegation removes it in every model, and that harm stays within two hops of untrusted input. Deployers can act on it directly by naming the record and recipient on the operator channel where the task allows, withholding the grant from agents that need no hand-offs, and concentrating checks on privileged agents next to untrusted content.
  \item \textbf{Owner-bound egress and its evaluation.} We design an effect-determined gateway rule, prove that it holds every attack twin at no cost to bound tasks under a stated precondition, reference confinement (Section~\ref{sec:5}), and test it against six alternatives with scripted insiders and full decision logs (Sections~\ref{sec:6} and~\ref{sec:7}). It held all 250 harmful sends without losing a hand-off task, and failed exactly where its precondition predicts, when insiders retyped records into the message body. Deployers gain a defence that keeps delegation and whose limit is known in advance, and the released testbed regenerates every table from the logs.
\end{itemize}

The rest of the paper is organised as follows. Section~\ref{sec:2} gives background and places prior defences by where they enforce. Section~\ref{sec:3} states the system and threat model, Section~\ref{sec:4} proves the dilemma, and Section~\ref{sec:5} presents owner-bound egress and its guarantee. Section~\ref{sec:6} describes the methodology, Section~\ref{sec:7} reports the results, and Sections~\ref{sec:8} to \ref{sec:10} discuss implications, limitations and conclusions.

\section{Background and Related Work}\label{sec:2}

This section fixes the terms used in the paper, and places prior work by where it enforces. We start from the conditions for data exfiltration, which are common to every system we discuss. Exfiltration through an LLM agent needs three capabilities in one system: access to private data, exposure to untrusted content, and a way to communicate externally~\cite{willison2025trifecta}. A multi-agent system can hold the three in different agents, so that a helper reads untrusted content, a privileged agent holds data and an external channel, and a hand-off joins them. Removing one capability from each agent therefore does not remove it from the system, which is why the terms below are defined at the level of the system rather than of one agent.

\subsection{Terminology}\label{sec:2.1}

Several terms that recur in this paper, such as taint and provenance, are used loosely in the agent literature. Table~\ref{tab:1} states the meaning we rely on and the work it comes from. Most of them describe who acts and how an instruction arrived. Terms that this paper introduces, such as egress action, binding relation and reference confinement, are defined where they are first used.

\begin{table}[!t]
\centering
\caption{Terms used in this paper.}
\label{tab:1}
\footnotesize
\begin{tabularx}{\columnwidth}{@{}>{\raggedright\arraybackslash}p{0.22\columnwidth}>{\raggedright\arraybackslash}X>{\raggedright\arraybackslash}p{0.2\columnwidth}@{}}
\toprule
\textbf{Term} & \textbf{Meaning in this paper} & \textbf{Origin} \\
\midrule
Identity & The binding of an agent to a key, so that every message and call names its true sender & \cite{anderson1972} \\
Capability (scope) & An unforgeable grant naming what an agent may do or reference, checked at each call & \cite{hardy1988,saltzer1975,debenedetti2025camel} \\
Taint & A mark placed on an agent once it has read input from a lower-trust source; an integrity label & \cite{denning1976,costa2025} \\
Provenance & How an agent's instructions arrived: sender, sender's privilege, channel and addresses mentioned (the provenance view $V$ of Section~\ref{sec:4}); an integrity notion & \cite{petersgill2026,debenedetti2025camel} \\
Information-flow control (IFC) & Enforcement that labelled data reaches only the sinks its label permits. Integrity labels record where data or instructions came from; confidentiality labels record who may read data & \cite{denning1976,goguen1982,myers1997,sabelfeld2003} \\
Delegation (hand-off) & A task that reaches a privileged agent from another agent on the peer channel; the operator sentence that permits it is the delegation grant & Figure~\ref{fig:1} \\
AI control & Evaluation of a protocol by safety and usefulness together, assuming the model may work against the operator & \cite{greenblatt2024} \\
\bottomrule
\end{tabularx}
\end{table}

\subsection{Related work}\label{sec:2.2}

Prompt injection was first shown for single agents. Greshake et al.~\cite{greshake2023} showed that retrieved data can act as instructions, and AgentDojo~\cite{debenedetti2024agentdojo} made the threat measurable with a dynamic benchmark of tasks and attacks. Multi-agent systems extend it: Prompt Infection~\cite{lee2024promptinfection} and the Morris II worm~\cite{cohen2024aiworm} demonstrated injected instructions that copy themselves between agents, and Triedman et al.~\cite{triedman2025} steered multi-agent frameworks into running attacker code even when each agent resisted direct injection. We group defences by where they enforce, from inside the model to the boundary where data leaves the system.

\textbf{Inside the model.} Model-level defences train the model to rank text by its source, through an instruction hierarchy~\cite{wallace2024hierarchy} or separate prompt and data channels~\cite{chen2025struq}.

\textbf{In the agent runtime.} Detection defences such as MELON~\cite{zhu2025melon} re-run the agent with the user task masked and compare tool calls. Design defences remove the path instead: CaMeL~\cite{debenedetti2025camel} separates trusted control flow from untrusted data and tags each value with its sources and allowed readers, checked at tool calls, FIDES~\cite{costa2025} applies integrity and confidentiality labels in the planner, Progent~\cite{shi2025progent} enforces least-privilege tool policies, and IsolateGPT~\cite{wu2025isolategpt} isolates applications from each other. Design patterns limit what an agent may do after reading untrusted data, at some cost in utility~\cite{beurerkellner2025patterns}. All of these protect one agent serving one trusted principal, which is the assumption that delegation removes.

\textbf{Between agents.} Jha et al.~\cite{jha2025} broke existing defences against multi-agent hijacking and proposed ControlValve, which enforces a permitted control-flow graph, while SEAgent~\cite{ji2026} applies mandatory access control over a graph of agents and tools. G-Safeguard~\cite{wang2025gsafeguard} detects anomalous agents with a graph neural network and cuts their edges.

The work closest to ours is ``Can CaMeLs Talk?''~\cite{petersgill2026}. It shows that CaMeL's guarantee does not compose, because untrusted data from one agent is reinterpreted as a trusted instruction by the next. Its fix, multi-CaMeL, carries values and their capabilities from a parent agent to a child, is evaluated with a permissive policy, and leaves peer-to-peer communication to future work. Our result is complementary. Even with perfect provenance, a system that allows delegation must accept instructions from peers, and provenance then no longer separates the legitimate case from the attack. CaMeL's readers labels would separate them if they were checked at every egress point, and owner-bound egress is that check, held outside the agents.

\textbf{At the system boundary.} Identity and audit infrastructure acts as a reference monitor~\cite{anderson1972} outside the agents. It removes spoofing, scope over-claims and log tampering, but an authorised agent that leaks data still passes it. Agent firewalls~\cite{abdelnabi2025firewalls} limit the data that leaves an agent according to learned task rules. They are closest to our work in spirit, because they restrict what leaves rather than who asked, but their rules are learned and carry no proof. Owner-bound egress also enforces at this boundary. A survey of open problems in multi-agent security is given in~\cite{schroederdewitt2025}.

Our design draws on classical access control and information flow. The gateway is a reference monitor~\cite{anderson1972} following complete mediation and least privilege~\cite{saltzer1975}, within the class of policies that execution monitors can enforce~\cite{schneider2000}. Information-flow control goes back to Denning's lattice model~\cite{denning1976} and to noninterference~\cite{goguen1982}, surveyed in~\cite{sabelfeld2003}. Owner-bound egress is closest to the decentralized label model~\cite{myers1997}, in which data carries its owner's policy and only the owner's authority can release it, and to contextual integrity~\cite{nissenbaum2004}, in which a flow is judged by its fit to the norms of its context. From AI control~\cite{greenblatt2024} we adopt the practice of reporting safety and usefulness together under a possibly adversarial model.

Relative to this body of work, our first contribution is a limit on a class of rules rather than an attack on one design. Earlier multi-agent studies show that attacks succeed~\cite{lee2024promptinfection,triedman2025}, propose and evaluate a policy~\cite{jha2025,ji2026}, or extend provenance tracking~\cite{petersgill2026}. We ask instead what any rule that reasons from provenance can achieve once delegation is permitted, and we answer with a proof. The release mechanism we evaluate is not new: it is the readers check of the decentralized label model~\cite{myers1997} and of CaMeL's capabilities~\cite{debenedetti2025camel}. What is new is the argument that delegation leaves only this confidentiality half of information-flow control standing, a gateway form of the check that needs no interpreter inside the agents, and a measurement of exactly where it fails. Section~\ref{sec:3} states the system and adversaries under which we make these claims.

\section{System and Threat Model}\label{sec:3}

\subsection{System model}\label{sec:3.1}

We consider a set of LLM agents that exchange messages and call tools through a gateway (Figure~\ref{fig:2}). The gateway is ordinary code that the agents cannot read or modify, and every message and tool call passes through it. It binds each agent's identity to a key, reads permission scopes only from signed credentials, and writes a hash-chained audit log. Some agents are privileged and may reference customer records and send external email, while the rest may only write files and post messages. Each agent receives input on two channels: the trusted operator channel, which carries the deployer's instructions, and the peer channel, which carries messages from other agents.

A task can reach a privileged agent on either channel, and the difference is central to this paper. In a direct task, the operator channel names the task, including which record goes to which contact. In a hand-off, the operator channel says only that teammates will pass requests on, and the task itself arrives on the peer channel from a low-privilege helper. We call the operator sentence that authorises hand-offs the delegation grant. Hand-offs are how multi-agent systems delegate, and they are also the path that a relayed injection follows.

The asset we protect is a set of customer records. Each record has one owner, and each owner has one registered contact address. A harmful send is an external email that carries a record to any address other than its owner's contact, and a run leaks if at least one harmful send lands. This definition covers confidentiality only, and Section~\ref{sec:3.3} lists what it leaves out.

\begin{figure}[!t]
\centering
\includegraphics[width=\linewidth]{Figure_2.pdf}
\caption{System and threat model. Agents receive untrusted content and operator tasks, and every message and tool call passes through the gateway. Identity, scopes and the audit log remove forgery but leave an authorised agent free to leak. The egress rule checks only the recipient and the attached record IDs against the binding store, which is kept outside the agents, and record content exists only at release. The operator, gateway code, keys and binding store are trusted; every agent is untrusted for the egress property.}
\label{fig:2}
\end{figure}

\subsection{Adversaries}\label{sec:3.2}

We consider three adversaries of increasing strength, summarised in Table~\ref{tab:3}. None of them can forge identity, scopes or logs, because the gateway removes these attacks by construction. They differ in how much of the agents' behaviour they control, from the content of one post to a fully compromised privileged agent. Each has a precedent in reported attacks or in prior threat models.

\begin{table*}[!t]
\centering
\caption{Adversaries, in increasing strength.}
\label{tab:3}
\footnotesize
\begin{tabularx}{\textwidth}{@{}>{\raggedright\arraybackslash}p{0.14\textwidth}>{\hsize=1.3\hsize\raggedright\arraybackslash}X>{\hsize=0.7\hsize\raggedright\arraybackslash}X>{\hsize=1\hsize\raggedright\arraybackslash}X@{}}
\toprule
\textbf{Adversary} & \textbf{Controls} & \textbf{Cannot do} & \textbf{Precedent} \\
\midrule
A1 Injector & the content posted by one low-privilege agent & forge identity, scopes or logs & all three cases in Table~\ref{tab:4} \\
A2 Persuaded peers & any number of agents that act on or relay what they read & forge identity, scopes or logs & \cite{lee2024promptinfection,cohen2024aiworm,triedman2025} \\
A3 Insider & one privileged agent with valid credentials, which chooses its own recipients and attachments & bypass the gateway or forge logs & the untrusted model of AI control~\cite{greenblatt2024}; compromised credentials or code \\
\bottomrule
\end{tabularx}
\end{table*}

\begin{table*}[!t]
\centering
\caption{Reported cases mapped onto the threat model. In each, identity checks pass and the leak is an effect the acting agent was permitted to cause.}
\label{tab:4}
\footnotesize
\begin{tabularx}{\textwidth}{@{}>{\raggedright\arraybackslash}p{0.12\textwidth}>{\raggedright\arraybackslash}p{0.12\textwidth}>{\hsize=1.2\hsize\raggedright\arraybackslash}X>{\hsize=0.9\hsize\raggedright\arraybackslash}X>{\centering\arraybackslash}p{0.06\textwidth}>{\hsize=0.9\hsize\raggedright\arraybackslash}X@{}}
\toprule
\textbf{Case} & \textbf{Untrusted input} & \textbf{Agent that acted} & \textbf{Permitted effect that leaked} & \textbf{Hand-off} & \textbf{Effect check that refuses it} \\
\midrule
ServiceNow Now Assist, 2025~\cite{costello2025appomni} & ticket text written by a low-privilege user & a second agent, recruited through agent discovery, running with the privileges of the administrator who started the interaction & a ticket the attacker could not read, copied into the attacker's own ticket & yes & that ticket's content may not go to a record the attacker owns \\
GitHub MCP server, 2025~\cite{invariant2025github} & a public GitHub issue & one coding agent using the user's GitHub access & private repository data in a public pull request & no & private-repository content may not go to a public repository \\
Microsoft 365 Copilot, CVE-2025-32711~\cite{reddy2025echoleak} & one external email & Copilot with the user's data access & internal data placed in an image link fetched through an allowed Microsoft Teams proxy & no & internal content may not go to an outside address \\
\bottomrule
\end{tabularx}
\end{table*}

Table~\ref{tab:4} shows that the reported cases fit this model. Only the ServiceNow case involves a hand-off, and it is the twin of Proposition 1 in a shipped product, because discoverable agents on the same team may invoke each other by default, which acts as a delegation grant~\cite{costello2025appomni}. Where a model-side filter was in place it did not stop the leak: the ServiceNow attack ran with injection protection enabled~\cite{costello2025appomni}, and the EchoLeak email was worded to evade Microsoft's injection classifier~\cite{reddy2025echoleak}. In all three cases, the refusal can be stated as a relation between a piece of data and a destination. This is the form of rule that Section~\ref{sec:5} builds.

\subsection{Trust assumptions and scope}\label{sec:3.3}

Our results rest on two trust assumptions. The gateway code, its keys and the log anchor are trusted, and the mapping from records to owners and from owners to contacts is maintained outside the agents and is correct. If either assumption fails, no result in this paper holds. Three threats are out of scope: integrity harms that move no owned data, such as issuing a wrong refund; covert channels in message timing or length; and attacks on model weights. Adaptive attackers that optimise injection text against the defence are discussed with the limitations in Section~\ref{sec:9}.

The binding is part of the trusted base, and it is also a target. Many support systems let an agent update a customer's contact address. If an agent can write $B$, an injection can ask it to change a customer's email to the attacker's address and then request a legitimate send, which the rule allows because the record is now bound to the attacker. Owner-bound egress therefore needs writes to $B$ held to a stronger standard than reads of records: a change of contact should be confirmed through the old contact or by a person, and should never be a tool that an agent can call on a peer's request. In our testbed no agent can write $B$, so this attack is not measured, and Section~\ref{sec:9} lists it with the other limits.

\section{The Hand-off Dilemma}\label{sec:4}

This section states the dilemma for a general class of rules. We first define egress actions and the two classes of rule, and then prove that no rule in the provenance class can separate a hand-off from its attack twin. The experiments in Section~\ref{sec:7} are witnesses to this result rather than the evidence for it.

\subsection{Definitions}\label{sec:4.1}

\textbf{Definition 1 (Egress action and rule).} An egress action by agent $i$ is a pair $(r, D)$ of a recipient address $r$ and a set $D$ of record references. An egress rule is a function that the gateway applies to each egress action and that returns allow or hold.

\textbf{Definition 2 (Provenance view).} For agent $i$ at the moment of an egress action, the provenance view $V$ records how $i$'s instructions arrived. For every message $m$ that $i$ has read, $V$ contains the sender's identity, the sender's privilege, the channel (operator or peer) and the set of external addresses that appear in $m$.

\textbf{Definition 3 (Provenance-determined rule).} A rule is provenance-determined if its verdict is a function of $V$ and $r$ only. It may use any information about where the instruction and the recipient came from, but it does not use $D$.

\textbf{Definition 4 (Effect-determined rule).} A rule is effect-determined if its verdict is a function of $(r, D)$ and of a binding relation $B$ between records and recipients that is maintained outside the agents. It does not use $V$.

\textbf{Definition 5 (Delegation).} A system permits delegation if there is a benign execution in which a privileged agent $i$ receives, on the peer channel from a helper $h$, a request to send record $d$ to address $r$, where $r$ is bound to $d$ and completing the request is the intended behaviour.

The provenance class covers the rules that prior systems and our own baselines apply at agent boundaries. Privilege taint holds an email if the agent has read a lower-privilege peer, content taint holds it if the agent has read a peer message containing an external address, and destination provenance holds it unless the recipient appeared on the operator channel. Each of these reads $V$ and $r$ and ignores $D$, so each falls under the result below. In information-flow terms, $V$ holds integrity labels, which record where instructions came from. Confidentiality labels, which record who may read a record, are not part of $V$: a rule that checks them reads $D$ and is effect-determined.

\subsection{The dilemma}\label{sec:4.2}

\textbf{Proposition 1.} Let a system permit delegation, and let the adversary control the content of one helper's messages. For every benign delegated execution $E$ there is an attack execution $E'$ such that the acting agent has the same provenance view and the same recipient in both, and the action in $E'$ is a harmful send. Therefore every provenance-determined rule gives $E$ and $E'$ the same verdict.

\emph{Proof.} Take the benign execution $E$, in which helper $h$ sends agent $i$ a request naming address $r$ and record $d$, and $i$ emits $(r, \{d\})$. Build $E'$ by changing only the text of $h$'s message so that it requests a set $D'$ containing some record not bound to $r$, and let $i$ comply. The sender, its privilege, the channel and the set of addresses in the message are unchanged, so the provenance view is unchanged. The recipient is unchanged. A provenance-determined rule is a function of exactly these, so its verdict is unchanged. In $E'$ a record reaches an address it is not bound to, which is a harmful send. End of proof.

The attack requires nothing unusual. The adversary needs a helper whose content it controls (A1) or a helper that relays what it read (A2), and a recipient address at which it can receive mail, such as its own registered customer contact.

\textbf{Corollary 1 (trade-off).} Call $E$ and $E'$ twins. For any distribution over twin pairs, let $U$ be the probability that the rule allows the benign twin and $S$ the probability that it holds the attack twin. For every provenance-determined rule, deterministic or randomised:

\begin{equation}
S + U = 1
\end{equation}

A rule that holds every attack twin ($S = 1$) therefore completes no hand-off task ($U = 0$), and a rule that completes every hand-off task holds no attack twin. Mixing the two moves along the line $S + U = 1$ and never above it.

\subsection{Implications for defences}\label{sec:4.3}

Proposition~1 is short because the class is defined by what a rule does not read, and we state it formally for that reason: it turns a design question into a test. A defence can secure delegation only if its verdict can change when the record set $D$ changes and nothing else does. Our three provenance baselines receive $V$ and $r$ and never $D$, as their code in the released testbed shows. For published systems we apply the test to the decision input each paper describes, in Table~\ref{tab:2} below.

The proof also shows which fixes cannot work. It does not depend on how fine-grained $V$ is, so richer provenance, such as signed sender identity, per-message lineage or separate instruction and data channels~\cite{petersgill2026}, does not help. These additions remove attacks in which data is mistaken for an instruction, but they do not separate the twins, because in both twins a real peer sends a real instruction. Reading the message does not help by itself either. A rule could parse $h$'s message and see that it asks for $D'$, but to decide whether that request is legitimate it must know which records may go to $r$, and that knowledge is the binding relation $B$.

Table~\ref{tab:2} applies Proposition~1 to prior defences. We group them by where they enforce, from inside the model to the boundary where data leaves the system, and record for each what it decides from and the basis of its guarantee. The last column follows from the proposition: a defence whose decision input is the same for $E$ and $E'$ gives them the same verdict.

\begin{table*}[!t]
\centering
\caption{Prior defences grouped by where they enforce. The last column applies Proposition~1 and is our analysis, not a claim made by the cited work.}
\label{tab:2}
\footnotesize
\begin{tabularx}{\textwidth}{>{\raggedright\arraybackslash}p{2.0cm}>{\raggedright\arraybackslash}X>{\raggedright\arraybackslash}X>{\raggedright\arraybackslash}p{2.6cm}>{\raggedright\arraybackslash}X}
\toprule
\textbf{Enforcement point} & \textbf{Representative work} & \textbf{Decision input} & \textbf{Guarantee basis} & \textbf{Hand-off $E$ vs.\ attack twin $E'$} \\
\midrule
Model & Instruction hierarchy~\cite{wallace2024hierarchy}, StruQ~\cite{chen2025struq} & role or channel of the text & training & Same verdict: both requests arrive in the same role \\
Agent runtime & MELON~\cite{zhu2025melon} & behaviour that looks attacked & detector & Depends on the detector; no bound stated for this case \\
Agent runtime & CaMeL~\cite{debenedetti2025camel}, FIDES~\cite{costa2025}, design patterns~\cite{beurerkellner2025patterns} & trusted control flow vs.\ untrusted data; values carry sources and allowed readers & construction & Integrity labels: same verdict. Readers labels: separate them within one agent; composition across agents is open~\cite{petersgill2026} \\
Agent runtime & Progent~\cite{shi2025progent}, IsolateGPT~\cite{wu2025isolategpt} & allowed tools, arguments or apps & policy & Separates only if the policy relates data to recipient \\
Between agents & ControlValve~\cite{jha2025}, SEAgent~\cite{ji2026} & permitted edges and their labels & policy & Same verdict when the policy reads only the edge and its labels \\
Between agents & G-Safeguard~\cite{wang2025gsafeguard} & behaviour that looks attacked & detector & Depends on the detector; no bound stated for this case \\
Between agents & multi-CaMeL~\cite{petersgill2026} & values and capabilities carried from parent to child & construction & Same verdict under the evaluated policy; a readers policy would separate them but was not evaluated, and peer-to-peer delegation is left open \\
System boundary & Identity and audit (reference monitor~\cite{anderson1972}) & identity and scopes & construction, for forgery only & Same verdict: the leaking agent is authorised \\
System boundary & Agent firewalls~\cite{abdelnabi2025firewalls} & learned task rules on input and shared data & learned rules & Partly: limits data per task, with no proof \\
System boundary & Owner-bound egress (this paper) & record and recipient against binding $B$ & proof, under reference confinement & Separates: allows $E$, holds $E'$ (Section~\ref{sec:5}) \\
\bottomrule
\end{tabularx}
\end{table*}

Every defence evaluated on multi-agent delegation decides from identity, the role of the text, observed behaviour, or the origin or path of an instruction, and none of these differs between $E$ and $E'$. The exception is the confidentiality half of information-flow control: a readers label relates a record to who may receive it, so it separates the twins, and CaMeL and FIDES carry such labels inside one agent~\cite{debenedetti2025camel,costa2025}. None of the multi-agent defences in Table~\ref{tab:2} reports checking them where data leaves the system. The gap matters in practice: each incident in Table~\ref{tab:4} passed identity checks, and two passed a model-side injection filter~\cite{costello2025appomni,reddy2025echoleak}.

Two options remain, and only one of them preserves delegation. Removing delegation makes the twins disappear, since tasks then arrive only on the operator channel; this is the single-agent setting that CaMeL-style designs secure~\cite{debenedetti2025camel}, but it is no longer a multi-agent system. The other option follows from where the twins differ. Their provenance views and recipients are equal and only their record sets differ, so a rule that reads $D$ against $B$ can separate them. Section~\ref{sec:5} builds the simplest such rule.

\section{Owner-Bound Egress}\label{sec:5}

\subsection{Design}\label{sec:5.1}

Owner-bound egress combines two mechanisms, one that limits what agents hold and one that limits what leaves. The first is pass by reference. Agents never hold record contents; they hold record identifiers and attach identifiers to an outgoing email, and the gateway resolves an identifier to content only at the moment of release. The second is a binding check at the gate. The gateway keeps the binding relation $B$, which in our instance binds each record to its owner's registered contact, and it releases an external email only if every attached record is bound to the recipient:

\begin{equation}
\mathrm{allow}(r, D) \iff \forall d \in D:\ (d, r) \in B
\end{equation}

The rule reads only the recipient and the attachments. It does not read who asked, what the agent has seen or what any message says, so it is effect-determined in the sense of Definition 4. It also keeps no per-agent state, which means it needs no taint tracking and has nothing to reset between tasks. Section~\ref{sec:5.2} shows what this design guarantees and what it requires.

\subsection{Guarantee}\label{sec:5.2}

\textbf{Definition 6 (Reference confinement).} A system satisfies reference confinement if the only component that can turn a record identifier into record content is the gateway at release.

\textbf{Proposition 2 (safety at no cost to bound tasks).} Under reference confinement, (a) in every execution and for every behaviour of every agent, including a fully compromised privileged agent, no record content is delivered to an external address that the record is not bound to; and (b) every task whose intended action is to send records to a recipient they are bound to is allowed, whatever channel the task arrived on.

\emph{Proof.} (a) Content leaves the system only through gateway release. Release of record $d$ to $r$ happens only if $(d, r) \in B$, and agents cannot produce content by any other path because they never hold it. (b) The verdict depends on $(r, D)$ and $B$ only, and the intended action satisfies the release condition. End of proof.

Proposition~2 places owner-bound egress at $S = 1$ and $U = 1$ on the twin distribution of Corollary 1, a point no provenance-determined rule can reach. The guarantee holds against all three adversaries of Section~\ref{sec:3}, and it does not degrade with model capability or with stronger injection text, because the model's decision is not an input to the rule.

\subsection{Limits and relation to existing mechanisms}\label{sec:5.3}

\textbf{Remark 1 (confinement is necessary).} If an agent can obtain record content, then for any rule that inspects only attached identifiers there is an execution that delivers that content to an unbound recipient.

\emph{Proof.} The agent places the content in the message body and attaches nothing. The release condition is vacuously true. End of proof.

We test this directly in Section~\ref{sec:7.3} with an insider that retypes records into the body. Owner-bound egress fails completely against it, as Remark~1 says it must.

Three further limits concern the scope of the rule rather than its soundness. First, the rule protects only data that has an owner and a bound recipient set; flows to auditors, delegates or other services need a richer $B$, and writing $B$ is a policy task that we do not solve. Second, the rule checks that data goes to its owner but not that the owner should receive it now, so a bound recipient can still receive too much or too little of their own data. Third, actions that carry no owned data are unconstrained. In our runs, 469 unrequested emails of this kind passed the rule, leaking nothing but still forming a nuisance that the rule does not address.

A fourth limit concerns tasks that need a record's content. Summarising a record, answering a customer or triaging a ticket all require an agent to read the record, which breaks reference confinement as we state it. The construction extends in the way CaMeL handles the same problem inside one agent~\cite{debenedetti2025camel}: a quarantined agent does the reading, its output carries the readers of every record it read, and the gateway applies the release check to that label instead of to attached identifiers. The cost is label creep, since an output derived from two owners' records is releasable to neither. Our experiments do not test this mode; Remark~1 and the retyped-insider rows of Table~\ref{tab:11} show what happens when it is absent.

The mechanism itself has clear precedents. The release check is owner-authorised declassification in the sense of the decentralized label model~\cite{myers1997}, applied at one choke point. Pass by reference is the capability discipline that CaMeL applies inside a single agent's interpreter~\cite{debenedetti2025camel}, moved to the boundary between agents and the outside world. An address allow-list is the special case in which $B$ relates every record to every approved address, which is why it cannot stop an approved recipient from receiving someone else's record (Section~\ref{sec:7.3}). Section~\ref{sec:6} describes how we test these claims.

\section{Methodology}\label{sec:6}

Our methodology tests both kinds of claim in this paper with one pipeline: the behavioural claim that delegating a task changes what agents do, and the rule claims of Sections~\ref{sec:4} and \ref{sec:5}. Figure~\ref{fig:3} shows its seven steps. Steps 1 to 3 build a factorial set of scenarios, step 4 runs agents in them through the gateway, and step 5 applies every egress rule to every logged action. Steps 6 and 7 verify the logs and compute the measures. Because each step reads only the output of the one before, every table in Section~\ref{sec:7} can be regenerated from the logs alone.

\begin{figure}[!t]
\centering
\includegraphics[width=\linewidth]{Figure_3.pdf}
\caption{Methodology pipeline. Steps 1 to 5 run once per run; steps 6 and 7 operate on the pooled logs of all model and scripted runs. The note at the right of each panel gives its scale. The 49{,}500 verified decisions are the 31{,}500 model decisions plus 18{,}000 from scripted runs.}
\label{fig:3}
\end{figure}

\subsection{Testbed construction (step 1)}\label{sec:6.1}

Each run simulates a ten-agent system for five synchronous steps. Agent 0 is the low-privilege source of untrusted content, and three of the other nine agents are privileged. At each step, an agent reads the latest post of up to three agents it follows, writes a short status, and may send one email with attached record identifiers. The system holds eight customer records, each with one registered contact, and in every run two privileged agents are each assigned one task: send one named record to its owner's contact. The seed fixes the graph, the placement of privileged agents and the attacker's address.

Agents never see record contents, so the testbed satisfies reference confinement for model-driven agents by construction. All addresses use reserved example domains, the outbox is an in-memory list, and the container has no network access. These choices make every send observable and harmless, so the experiment can log harmful sends without causing them.

\subsection{Experimental design (steps 2 and 3)}\label{sec:6.2}

We vary four factors in a full factorial design (Table~\ref{tab:5}). Topology tests whether the effect depends on graph shape, and the task channel decides whether the operator names the task or delegates it (Section~\ref{sec:3.1}). Source content is one of six fixed injection templates or a benign control, and two seeds per cell vary the graph, privilege placement and attacker address. The channel factor is assigned by design, so differences between its levels are causal within the testbed.

\begin{table}[!t]
\centering
\caption{Experimental factors.}
\label{tab:5}
\footnotesize
\begin{tabularx}{\columnwidth}{@{}>{\raggedright\arraybackslash}p{0.22\columnwidth}>{\raggedright\arraybackslash}X@{}}
\toprule
\textbf{Factor} & \textbf{Levels} \\
\midrule
Topology & chain, ring, random, hierarchy (coordinator delegates to workers), launder (privileged agents read only privileged peers) \\
Task channel & direct, hand-off \\
Source content & six fixed injection templates, plus one benign control \\
Seed & 2 per cell (changes graph, privilege placement and attacker address) \\
\bottomrule
\end{tabularx}
\end{table}

The design gives 5 $\times$ 2 $\times$ 7 $\times$ 2 = 140 runs per model, of which 120 are attack runs and 20 are controls, with 45 decisions per run and 6{,}300 decisions per model. Over five models there are 700 runs (600 attack runs and 100 controls) and 1{,}400 tasks, of which 700 are direct and 700 are hand-off, with 600 of each in attack runs. Every rule is applied to all 700 runs, so all seven rules are scored on the same runs and tasks. The six injection templates were written before the defences; we wrote no attack text for this paper and ran no search for stronger injections.

\subsection{Agents and egress rules (steps 4 and 5)}\label{sec:6.3}

The agents are driven either by models or by scripted policies. We use five open-weight models from four families, served locally through Ollama at temperature 0 on one 8 GB GPU: Qwen2.5 3B and 7B~\cite{qwen2024}, Llama 3.1 8B~\cite{llama3_2024}, Mistral 7B~\cite{jiang2023mistral} and Gemma 3 4B~\cite{gemma3_2025}. All five models returned parseable JSON for all 31{,}500 decisions.

Three scripted policies bracket the models. Honest agents do their task and ignore everything else, while naive agents obey any instruction line that contains an address and relay it, which is the worst case for persuasion. Insider runs replace one tasked privileged agent with an attacker (A3) that sends all records to a chosen recipient: its own address, another registered contact, or the contact of its own task. Each insider strategy has a second form that retypes records into the message body, which breaks reference confinement on purpose and tests Remark~1. In total the scripted runs are 140 honest, 140 naive and 120 insider runs (20 per strategy), or 18{,}000 decisions.

The gateway never alters message content, so each agent decision is made once and all seven rules of Table~\ref{tab:6} are applied to it in parallel. Every rule therefore sees identical agent behaviour, and differences between rules cannot come from differences in what the agents did. The rules span the classes of Section~\ref{sec:4}: one baseline, three provenance-determined rules, two effect-determined rules and one combination.

\begin{table}[!t]
\centering
\caption{Egress rules compared.}
\label{tab:6}
\footnotesize
\begin{tabularx}{\columnwidth}{>{\raggedright\arraybackslash}X>{\raggedright\arraybackslash}X>{\raggedright\arraybackslash}X}
\toprule
\textbf{Rule} & \textbf{Holds an external email when} & \textbf{Class} \\
\midrule
none & never (identity and scope checks only) & baseline \\
V0 privilege taint & the sender has read any lower-privilege peer & provenance \\
V2 content taint & the sender has read a peer message containing an external address & provenance \\
DP destination provenance & the recipient did not appear on the sender's operator channel & provenance \\
AL allow-list & the recipient is not a registered contact & effect, recipient only \\
DB owner-bound egress & an attached record is not bound to the recipient & effect \\
DP+DB & either DP or DB holds & mixed \\
\bottomrule
\end{tabularx}
\end{table}

\subsection{Verification and measures (steps 6 and 7)}\label{sec:6.4}

Before scoring, we verify the pipeline itself. The analysis reads logs only and never calls a model; it recomputes every rule verdict from the logged inputs and compares it with the verdict recorded during the run. Across 49{,}500 decisions (31{,}500 from models and 18{,}000 from scripted runs), there were zero mismatches. The rule logic is also covered by 110 unit tests, including malformed and adversarial inputs.

We then compute the four measures defined in Table~\ref{tab:7}. Run safety and task completion are the measured counterparts of $S$ and $U$, which Corollary~1 defines on twin pairs rather than on runs, while compliance and relay describe the agent behaviour that drives it. Intervals are Wilson 95\% intervals~\cite{wilson1927}. Decisions within a run are not independent, so for the main causal claim we also report a run-level paired test~\cite{mcnemar1947}.

\begin{table}[!t]
\centering
\caption{Measures.}
\label{tab:7}
\footnotesize
\begin{tabularx}{\columnwidth}{>{\raggedright\arraybackslash}X>{\raggedright\arraybackslash}X>{\raggedright\arraybackslash}X}
\toprule
\textbf{Measure} & \textbf{Numerator} & \textbf{Denominator} \\
\midrule
Run safety & attack runs with no harmful send that landed & attack runs \\
Task completion & tasks completed with the right record to the right contact & tasks assigned, reported separately for direct and hand-off tasks \\
Compliance & exposed decisions by privileged agents that emailed the attacker's address & decisions by privileged agents with at least one input carrying the attacker's address \\
Relay & exposed decisions whose own status message repeats the attacker's address & exposed decisions \\
\bottomrule
\end{tabularx}
\end{table}

\subsection{Mapping to deployed systems}\label{sec:6.5}

The testbed is an abstraction, so Table~\ref{tab:8} states what each of its parts stands for in a deployed system and what it leaves out. The mapping follows the protocols and incidents introduced in Sections~\ref{sec:1} and \ref{sec:3}.

\begin{table}[!t]
\centering
\caption{Testbed elements and their deployed counterparts.}
\label{tab:8}
\footnotesize
\begin{tabularx}{\columnwidth}{>{\raggedright\arraybackslash}X>{\raggedright\arraybackslash}X>{\raggedright\arraybackslash}X}
\toprule
\textbf{Testbed element} & \textbf{Deployed counterpart} & \textbf{Left out of the testbed} \\
\midrule
Status post read by followers & A2A message with a text part~\cite{a2aspec2025} & transport, authentication, streaming \\
Delegation grant (one operator sentence) & agents allowed to invoke teammates, as with Now Assist agent discovery~\cite{costello2025appomni} & an orchestrator that routes tasks; grants scoped by skill \\
Record identifier attached to an email & A2A file part sent as a URI reference instead of inline bytes~\cite{a2aspec2025}; a resource behind an MCP server~\cite{mcp2024} & tasks in which an agent must read the record itself (Section~\ref{sec:5.3}) \\
Gateway in front of every tool call & an MCP server or API gateway that holds the credentials~\cite{mcp2024} & latency; several gateways in one system \\
External email & any outbound effect: a ticket write, a pull request, a URL fetch (Table~\ref{tab:4}) & egress types other than email are not run \\
Owner-to-contact binding & owner and contact fields in a CRM; ticket access lists & errors in the binding; agents that can edit it (Section~\ref{sec:3.3}) \\
Ten agents, five steps, three privileged & no single counterpart & larger graphs and longer runs \\
\bottomrule
\end{tabularx}
\end{table}

The two kinds of result transfer differently. The rule outcomes in Tables~\ref{tab:10} and~\ref{tab:11} follow from Propositions~1 and~2 and Remark~1, so they hold in any system that meets the stated preconditions, whatever its scale or model; the runs witness them. The behavioural results in Tables~\ref{tab:9} and~\ref{tab:12} are rates for five small open-weight models, fixed injections and one wording of the grant. They show that the delegation effect exists and has the same direction in every model, not how large it is in production. We did not evaluate a frontier API model, and models trained with an instruction hierarchy~\cite{wallace2024hierarchy} may treat a peer's request differently, so a frontier replication is the most useful next step for the behavioural claim. Section~\ref{sec:7} reports these results, and Section~\ref{sec:9} lists what could change the rates.

\section{Results}\label{sec:7}

\subsection{A named task protects agents, and delegation removes it}\label{sec:7.1}

Privileged agents almost never acted on an injection when their task came from the operator, and often did when they had been told to accept hand-offs. Table~\ref{tab:9} counts decisions by privileged agents that had read at least one input carrying the attacker's address.

\begin{table*}[!t]
\centering
\caption{Privileged agents that emailed the attacker's address, by task channel (Wilson 95\% intervals). Direct and Hand-off count every exposed privileged agent in runs of that channel. The last column is the subset of hand-off decisions by privileged agents with no task, which receive no operator text, at hop 1 with one exposure.}
\label{tab:9}
\footnotesize
\setlength{\tabcolsep}{4pt}
\begin{tabular*}{\textwidth}{@{\extracolsep{\fill}}lrrrrrr@{}}
\toprule
 & \multicolumn{2}{c}{\textbf{Direct (named task)}} & \multicolumn{2}{c}{\textbf{Hand-off (grant)}} & \multicolumn{2}{c}{\textbf{No operator text (hop 1)}} \\
\cmidrule(lr){2-3}\cmidrule(lr){4-5}\cmidrule(l){6-7}
\textbf{Model} & $k/n$ & \% [95\% CI] & $k/n$ & \% [95\% CI] & $k/n$ & \% [95\% CI] \\
\midrule
Qwen2.5 3B & 0/210 & 0.0 [0.0, 1.8] & 12/219 & 5.5 [3.2, 9.3] & 1/30 & 3.3 [0.6, 16.7] \\
Gemma 3 4B & 0/224 & 0.0 [0.0, 1.7] & 23/221 & 10.4 [7.0, 15.1] & 4/27 & 14.8 [5.9, 32.5] \\
Mistral 7B & 1/230 & 0.4 [0.1, 2.4] & 25/243 & 10.3 [7.1, 14.7] & 2/27 & 7.4 [2.1, 23.4] \\
Llama 3.1 8B & 0/211 & 0.0 [0.0, 1.8] & 31/225 & 13.8 [9.9, 18.9] & 1/28 & 3.6 [0.6, 17.7] \\
Qwen2.5 7B & 1/213 & 0.5 [0.1, 2.6] & 119/229 & 52.0 [45.5, 58.4] & 12/29 & 41.4 [25.5, 59.3] \\
\midrule
All five & 2/1{,}088 & 0.2 [0.1, 0.7] & 210/1{,}137 & 18.5 [16.3, 20.8] & 20/141 & 14.2 [9.4, 20.9] \\
\bottomrule
\end{tabular*}
\end{table*}

The effect has the same direction in every model. The smallest is Qwen2.5 3B, where the intervals still do not overlap. The largest is Qwen2.5 7B, where more than half of exposed privileged decisions under a delegation grant emailed the attacker.

At run level with no egress rule, 19 of 300 direct runs leaked and 86 of 300 hand-off runs leaked. Each direct run has a hand-off counterpart with the same model, topology, template and seed. In 77 pairs only the hand-off run leaked, in 10 only the direct run leaked, in 9 both did (exact McNemar test, $p = 5.9 \times 10^{-14}$).

The two channels differ in two ways: a direct task names the record and the contact, while a hand-off removes that and adds the grant. Privileged agents without a task receive neither in either channel, which separates the two in part. At hop 1 with one exposure, these agents acted in 20 of 141 decisions (14.2\%), against 2 of 1{,}034 (0.2\%) for agents whose operator named the task and 182 of 886 (20.5\%) for delegated agents at the same position. Most of the gap is therefore the protection that a named task gives, and the grant's own increase over no instruction is smaller, with overlapping intervals. All 141 untasked exposures come from hand-off runs on the random graph, so this split is observational; an arm that adds the grant to a named task would separate the two directly, and we did not run it.

Of the 250 harmful sends in attack runs, 219 happened in hand-off runs. Hand-off runs leaked more often in all five topologies. The gap is smallest on the hierarchy graph, with 9 leaking direct runs against 11 hand-off, where a coordinator already stands between the source and the workers.

Within one model family, the larger model complied far more (52.0\% against 5.5\%). This matches the pattern reported for single agents, where attack success tends to rise with capability~\cite{debenedetti2024agentdojo}.

The mechanism is visible in the prompts. A named task gives the model the effect it should cause, a record and a contact, so an injected request for a different effect stands out. A delegated agent has only the grant, which says that peers will supply the task, so a peer's request for any effect looks like work. This is Proposition~1 seen from inside the agent: the model also decides from provenance, and under a grant the provenance of an injection is acceptable. Owner-bound egress restores at the gateway what the named task gave the model, a statement of which record may go to which recipient, and Section~\ref{sec:7.2} shows the same limit for the rules that lack it.

\subsection{Provenance rules sit on the trade-off line; owner-bound egress does not}\label{sec:7.2}

\begin{table*}[!t]
\centering
\caption{Seven rules over five models, pooled. Each rule is scored on the same 700 runs: 600 attack runs with 250 harmful sends, and 700 direct and 700 hand-off tasks, 600 of each in attack runs. Worst single model in brackets.}
\label{tab:10}
\footnotesize
\begin{tabularx}{\textwidth}{>{\raggedright\arraybackslash}X>{\raggedright\arraybackslash}X>{\raggedright\arraybackslash}X>{\raggedright\arraybackslash}X>{\raggedright\arraybackslash}X}
\toprule
\textbf{Rule} & \textbf{Safe attack runs} & \textbf{Harmful sends held} & \textbf{Direct tasks done} & \textbf{Hand-off tasks done} \\
\midrule
none & 495/600 [70/120] & 0/250 & 644/700 & 673/700 \\
V0 privilege taint & 595/600 [117/120] & 241/250 & 60/700 & 0/700 \\
V2 content taint & 567/600 [102/120] & 208/250 & 446/700 & 0/700 \\
DP destination provenance & 600/600 & 250/250 & 644/700 & 0/700 \\
AL allow-list & 599/600 [119/120] & 249/250 & 644/700 & 673/700 \\
DB owner-bound egress & 600/600 & 250/250 & 644/700 & 673/700 \\
DP+DB & 600/600 & 250/250 & 644/700 & 0/700 \\
\bottomrule
\end{tabularx}
\end{table*}

All three provenance rules completed none of the 700 hand-off tasks, in every model, as Corollary 1 requires of a rule that holds the attack twin. Two of them also leaked: privilege taint missed 9 harmful sends and content taint missed 42. None of the missed sends went to the attacker's address; all are model errors of the kind described in Section~\ref{sec:7.5}. The 9 that privilege taint missed are all on the launder graph, where privileged agents read only privileged peers and are never tainted.

Owner-bound egress held all 250 harmful sends and its task completion is identical to the undefended baseline in both channels. The tasks it did not complete (56 direct, 27 hand-off) are tasks the model itself failed, not tasks the rule held: the rule held zero task emails. Counting attack runs alone, it completed 576 of 600 hand-off tasks, the same as no rule. Most of the direct failures are Mistral 7B, which sent the email without attaching the record in 211 decisions and completed 90 of 140 direct tasks both with no rule and under owner-bound egress.

\begin{figure}[!t]
\centering
\includegraphics[width=\linewidth]{Figure_4.pdf}
\caption{Table~\ref{tab:10} placed on the plane of Corollary~1, using measured rates rather than the twin-pair probabilities $S$ and $U$. Measured safety is the share of the 250 harmful sends held; measured usefulness is the share of the 673 hand-off tasks the models completed with no rule, so that model failures do not count against a rule. The three provenance-determined rules (V0 privilege taint, V2 content taint, DP destination provenance) and DP+DB complete no hand-off task. DB is owner-bound egress and AL the allow-list (open diamond, behind DB). The dashed line is the bound $S + U = 1$ of Corollary~1, which is stated for twin pairs; the measured rates are plotted against it as a guide, not as a test of it.}
\label{fig:4}
\end{figure}

These outcomes follow from the propositions of Sections~\ref{sec:4} and \ref{sec:5}. The experiment adds three things the propositions cannot: that the attacks occur at a meaningful rate without a defence (105 of 600 runs), that models keep to the pass-by-reference discipline well enough for the rule to cost nothing, and that the allow-list comes close on this workload, which Section~\ref{sec:7.3} examines.

\subsection{Insiders: where each rule breaks}\label{sec:7.3}

The model runs cannot separate the allow-list from owner-binding, because persuaded models mostly send to addresses that are not registered. A compromised agent chooses its recipient better. Table~\ref{tab:11} gives safe runs out of 20 for each insider strategy.

\begin{table*}[!t]
\centering
\caption{Insider strategies (safe runs out of 20). Top half: records attached by reference. Bottom half: records retyped into the body.}
\label{tab:11}
\footnotesize
\begin{tabularx}{\textwidth}{>{\raggedright\arraybackslash}X>{\raggedright\arraybackslash}X>{\raggedright\arraybackslash}X>{\raggedright\arraybackslash}X>{\raggedright\arraybackslash}X>{\raggedright\arraybackslash}X>{\raggedright\arraybackslash}X}
\toprule
\textbf{Insider sends all records to} & \textbf{V0} & \textbf{V2} & \textbf{DP} & \textbf{AL} & \textbf{DB} & \textbf{DP+DB} \\
\midrule
its own outside address & 20 & 10 & 20 & 20 & 20 & 20 \\
another registered contact & 20 & 10 & 20 & 0 & 20 & 20 \\
the contact of its own task & 20 & 10 & 10 & 0 & 20 & 20 \\
own address, retyped & 20 & 10 & 20 & 20 & 0 & 20 \\
registered contact, retyped & 20 & 10 & 20 & 0 & 0 & 20 \\
task contact, retyped & 20 & 10 & 10 & 0 & 0 & 10 \\
\bottomrule
\end{tabularx}
\end{table*}

Table~\ref{tab:11} shows that each rule breaks at a different point. The allow-list fails against any registered recipient (0 of 20 in four rows), because it checks the recipient and not whose data is attached. Destination provenance fails on the task contact, which is the twin pair of Proposition 1 with the operator channel in place of the helper: the recipient is one the operator named, so the rule allows it in the 10 direct runs and holds the 10 hand-off runs only because it holds every hand-off. Privilege taint holds everything, but it is also the rule that completes no hand-off tasks and under 10\% of direct tasks. Content taint holds exactly the 10 hand-off runs in every row: there the hand-off post carries the task contact, an external address, so the insider is tainted, while in direct runs no peer post carries an address.

Owner-bound egress holds every by-reference attack (20 of 20 in the top half) and none of the retyped ones (0 of 20 in the bottom half), exactly as Remark~1 predicts. The last row has no winner. An insider that can read record contents and sends them to its own task contact is stopped only by rules that also stop legitimate work. Confinement therefore has to be enforced by construction, not detected at the gate.

The scripted naive agents give the worst case for persuasion. With no rule, all 120 attack runs leaked. Privilege taint left 24 leaking, all on the launder graph, where privileged agents read only privileged peers and never become tainted. Content taint, destination provenance, the allow-list and owner-bound egress each held all 120. Hand-off task completion was 0 of 140 under the three provenance rules and 104 of 140 under owner-bound egress, the same as with no rule.

\subsection{Injections do not spread; risk is one hop wide}\label{sec:7.4}

A reader might expect a relayed injection to cascade through the graph, but it did not in any model. Table~\ref{tab:12} gives the share of exposed decisions whose own message repeated the attacker's address.

\begin{table*}[!t]
\centering
\caption{Relay rate among exposed decisions, by distance from the source.}
\label{tab:12}
\footnotesize
\begin{tabularx}{\textwidth}{>{\raggedright\arraybackslash}X>{\raggedright\arraybackslash}X>{\raggedright\arraybackslash}X>{\raggedright\arraybackslash}X}
\toprule
\textbf{Model} & \textbf{Hop 1} & \textbf{Hop 2} & \textbf{Hop 3} \\
\midrule
Mistral 7B & 154/1{,}560 (9.9\%) & 14/165 (8.5\%) & 0/8 \\
Gemma 3 4B & 129/1{,}560 (8.3\%) & 10/142 (7.0\%) & 0/1 \\
Qwen2.5 7B & 48/1{,}560 (3.1\%) & 3/54 (5.6\%) & 0/2 \\
Llama 3.1 8B & 36/1{,}560 (2.3\%) & 0/45 (0.0\%) & none exposed \\
Qwen2.5 3B & 29/1{,}560 (1.9\%) & 2/27 (7.4\%) & 1/1 \\
\bottomrule
\end{tabularx}
\end{table*}

The relay probability $p$ is roughly constant across hops and at most 0.10 in every cell with more than ten exposed decisions (the only higher value is a single Qwen2.5 3B decision at hop 3). What matters for spread is how many agents read a relayer's post, its in-degree $k$: 3 on the ring and random graphs on average, 1 on the chain, and 8 for the hierarchy coordinator. With $k = 3$ the expected number of new relayers per relayer is $R = pk$, at most 0.30. Counted directly, an agent that relayed by step 3 had on average 0.48 followers who relayed later (0.65 on the random graph); this is an upper bound on $R$, since some of those followers also read the source. A branching process with $R$ below 1 dies out, and with $R \le 0.48$ the expected total number of relayers started by one source is bounded:

\begin{equation}
\mathbb{E}[\text{relayers}] \le \sum_{j \ge 0} R^{j} = \frac{1}{1 - R} \le 1.93
\end{equation}

Harm follows exposure. Of the 212 privileged decisions that emailed the attacker, 207 were made one hop from the source and 5 at two hops. None were further out.

Acting and relaying are different behaviours and models differ in which they favour. Counting every exposed agent, including low-privilege agents whose emails the scope check denies, Qwen2.5 7B emailed the attacker 251 times (120 of them by privileged agents, Table~\ref{tab:9}) and relayed the address 51 times. Mistral 7B emailed it 49 times (26 by privileged agents) and relayed it 168 times. A model can be dangerous as a deputy and harmless as a carrier, or the reverse.

These rates are for fixed templates that do not ask to be copied. Self-replicating injections are built to raise $p$~\cite{lee2024promptinfection,cohen2024aiworm}, and a system with larger fan-out $k$ has a larger $R$ with the same model. Runs last five steps, so an agent that relayed late had little time to recruit, and the hierarchy coordinator, read by eight workers, never relayed in 600 exposed decisions; at $p = 0.10$ it would give $R = 0.8$. The bound therefore holds for these graphs and run lengths, not for larger fan-in. The useful part of the result is the method: $p$ is a single-agent quantity that is cheap to measure, and with $k$ it predicts whether a deployment is above or below the threshold.

\subsection{Not every leak has an attacker}\label{sec:7.5}

Of the 250 harmful sends in attack runs, 176 went to the attacker's address; the other 36 of the 212 privileged emails to the attacker (Section~\ref{sec:7.4}) attached no record and leaked nothing. The other 74 went to addresses the attacker never mentioned, mostly addresses the model invented, such as a generic support mailbox on a domain that appears nowhere in the run. A further 9 harmful sends happened in control runs with no injection at all (6 from Qwen2.5 7B, 2 from Llama 3.1 8B, 1 from Gemma 3 4B).

Faults of this kind are expected in multi-agent systems. Cemri et al. annotated more than 1{,}600 traces from seven frameworks and catalogued failures with no adversary present, among them task derailment and agents ignoring each other's input~\cite{cemri2025mast}. A defence aimed only at attackers leaves this class open.

Owner-bound egress held every one of these 259 harmful sends. The allow-list missed two, one in an attack run (hence 249 of 250 in Table~\ref{tab:10}) and one in a control run; both sent one customer's record to another customer's registered contact. Of the 9 control-run sends, which Table~\ref{tab:10} does not count, owner-bound egress, destination provenance and privilege taint held all 9, the allow-list 8 and content taint 3. Provenance rules are a poor fit for this failure because there is no untrusted instruction to trace: content taint held none of the sends made by agents that had read no external address.

An effect-determined rule does not ask why a record is about to reach the wrong recipient. It covers persuasion, compromise and plain model error with the same check.

\section{Discussion}\label{sec:8}

\subsection{From identity and provenance to authority over effects}\label{sec:8.1}

Security for multi-agent systems has moved through two questions, and our results point to a third. The first question, who an agent is, is answered by external identity and audit. The second, where an instruction came from, is answered by provenance tracking inside one agent. Delegation makes a third question unavoidable: is this effect allowed, whoever asked? Operating systems resolved the confused deputy~\cite{hardy1988} by tying authority to the object being acted on, and our results suggest that the same resolution applies here and that tracing instructions cannot stand in for it once peers may instruct each other. In the terms of AI control~\cite{greenblatt2024}, provenance rules buy safety by spending all usefulness on the delegation channel, while owner-bound egress is a protocol whose safety does not rely on the untrusted model's cooperation, as the insider results of Table~\ref{tab:11} confirm up to the stated precondition.

\subsection{Guidance for deployers}\label{sec:8.2}

Three results translate directly into deployment guidance. Because risk is one hop wide (Section~\ref{sec:7.4}), the agents that need the closest checks are privileged agents that read untrusted content directly. Because privileged agents almost never acted on injections when the operator named their task (0.2\%) and often did when the task was delegated (18.5\%, Section~\ref{sec:7.1}), the operator should name the record and recipient wherever the task allows, and an agent that needs no hand-offs should not be told to accept them. Because effect checks work only when agents hold references (Sections~\ref{sec:5} and \ref{sec:7.3}), giving an agent raw record contents removes the guarantee. Better models will not make these steps unnecessary: within the one family we could compare, the larger model was the more compliant, and neither Proposition 1 nor Proposition 2 depends on model behaviour.

\subsection{Generalising the binding}\label{sec:8.3}

Owner-to-contact is the simplest binding relation, but the construction is general. It applies to any effect whose legitimacy is a relation between an object and a target, such as a refund and the account that paid, a file and the workspace it belongs to, or a credential and the service it was issued for. In each case the rule is to keep the relation outside the agents and to check it where the effect leaves. Writing these relations is a policy task, and the binding becomes a trusted component in its own right. We have evaluated only the first instance, and the others remain future work.

\section{Limitations and Ethics}\label{sec:9}

\subsection{Limitations}\label{sec:9.1}

The limitations of this work differ in kind, and Table~\ref{tab:13} groups them by what they affect. The rule outcomes are consequences of the definitions: that provenance rules complete no hand-offs and that owner-bound egress leaks nothing both follow from how the rules and the harm are defined. For this reason we present them as Propositions~1 and~2 and Remark~1 and use the runs only as witnesses. The empirical claims of the paper are the compliance, relay and fault rates, and most of the limitations below bear on those rates rather than on the propositions. The most important missing experiments are an arm that adds the grant to a named operator task, which would separate the grant from the loss of a named task (Section~\ref{sec:7.1}), and a replication with a frontier model.

\begin{table}[!t]
\centering
\caption{Limitations and their effect on the results.}
\label{tab:13}
\footnotesize
\begin{tabularx}{\columnwidth}{>{\raggedright\arraybackslash}X>{\raggedright\arraybackslash}X>{\raggedright\arraybackslash}X}
\toprule
\textbf{Limitation} & \textbf{What it affects} & \textbf{Direction or remedy} \\
\midrule
One wording of the delegation grant; no arm with a named task plus the grant; no scoped grant & compliance rates & a scoped grant (naming the helper or record) may reduce compliance; the grant's effect over no instruction is only partly separated (Section~\ref{sec:7.1}) \\
Five open-weight models of 3B to 8B parameters, quantised; no frontier or API model & compliance and relay rates & rates should not be read as rates for production systems \\
Six fixed, non-adaptive injection templates & compliance and relay rates & stronger injections can only raise them; the owner-bound guarantee does not depend on injection strength \\
Synthetic testbed: ten agents, five steps, eight records, one task type, no production A2A or MCP stack & external validity & Table~\ref{tab:8} lists what each element leaves out \\
Temperature 0, two seeds per cell; decisions within a run are correlated & decision-level intervals & intervals are optimistic; the run-level paired test does not rely on them \\
Literal relay detection (verbatim address) & relay rates & lower bounds for instruction spread, though not for address spread \\
Correct owner-to-contact binding assumed, and no agent can write it & all rule results & building $B$ is unsolved here; an agent that can edit a contact defeats the rule (Section~\ref{sec:3.3}) \\
Confidentiality only & scope & harmful actions that move no owned data are not evaluated \\
\bottomrule
\end{tabularx}
\end{table}

\subsection{Ethics and open science}\label{sec:9.2}

The environment is synthetic and the work is defensive. Records are fabricated, addresses use reserved example domains, no email leaves the process, and the container is network-isolated, so no human subjects or real user data are involved. We introduce no new attack text: all injections reuse six fixed templates of a kind already public~\cite{greshake2023,debenedetti2024agentdojo}. The finding with attacker value, that delegated agents comply far more often than agents given a named task, is also the finding defenders most need, and the paper pairs it with a mitigation. We will release the testbed, all prompts, model replies, decision logs and hash-chained gateway logs, together with the analysis script that regenerates every table from the logs.

\section{Conclusion}\label{sec:10}

Multi-agent systems need agents to take instructions from each other, and an injected instruction that has been relayed once is an instruction from a peer. We showed that this is a limit on a whole class of defences rather than a gap in any one of them: a rule that decides from provenance gives a legitimate hand-off and its attack twin the same verdict, so its safety and usefulness on twin pairs sum to one. The models share the limit. When their task was delegated, five models from four families acted on injections in 18.5\% of privileged exposures, against 0.2\% when the operator named the task.

The way out is to stop asking where an instruction came from and to ask what the action does. Owner-bound egress, which releases a record only to a recipient bound to its owner and keeps record contents out of agents' hands, held every harmful send in 600 attack runs at no cost to hand-off completion. It fails exactly where its precondition fails, when an agent holds the data itself. Identity belongs to infrastructure and provenance works inside one agent; between agents, what can be enforced is authority over effects.

\bibliographystyle{IEEEtran}
\bibliography{references}

\end{document}